\section{Crystalline cohomology with bounded divided powers}

% BEGIN D72P1
\paragraph{D72P1. The bar complex on a smooth affine lift.}
Let $R\twoheadrightarrow R_0=R/I$ carry the prescribed divided
powers, with $I^{[b]}=0$ for some $b\ge1$. Here $I^{[b]}$ is
generated by products of divided powers of total order at least $b$.
This is exactly the list condition in the source. Consider a smooth
affine $U=\operatorname{Spec}A$ over $R$ and its reduction $U_0$.
A test object is a surjection $B\to B_0$ with kernel $J$, compatible
divided powers with $J^{[m]}=0$ for some $m$, and a point of $U_0(B_0)$.
The coefficient at this object is $B$. In the ring direction both the
coefficient and the functor $Y$ of lifts of the given point to $U(B)$
are covariant. The set $Y(B)$ is nonempty: $J^m=0$ and $A$ is formally
smooth. For the site of $U$ over $(R,0)$ use its given $U(B_0)$-point
instead. The same assertion holds.

The cohomology of the category bar complex with this coefficient is
the derived inverse limit. Indeed the bar resolution of the constant
free rank-one diagram is a resolution by sums of free representable
diagrams, split after evaluation at any object by inserting its identity.
In degrees one and two its differential is exactly the alternating
composition and coefficient transport of the source. Another
resolution of the same constant diagram is the augmented complex
freely generated by ordered tuples $Y^{q+1}$: at each test it is the
nerve of the nonempty set of lifts, contractible by adjoining any one
of them.

In degree $q$ this latter functor of tuples is the filtered union
(colimit) of the representables of the divided-power diagonal envelopes
$D_{q,m}$. Form the relative divided-power envelope for the immersion
$U_0\hookrightarrow U^{q+1}$, kill divided-power order $m$, and use
$U_0$ as its reduction. The universal property is the choice of $q+1$
lifts over a single reduced point. Every test has some order bound,
so every tuple is represented. More precisely the universal map from an
envelope is unique once the tuple is fixed, and two bounds can be
increased to their maximum. This identifies the colimit as a functor.
The absolute construction uses
$U\hookrightarrow U^{q+1}$ and reduction $U$. Free representables
have higher Ext zero against $B\mapsto B$. For their sequential
colimit the telescope resolution calculates higher Ext by the
two-term $1-\mathrm{shift}$ on the product of the $D_{q,m}$.
Each transition is a quotient. Thus this difference map is onto
(solve one coordinate after the preceding one), giving Ext zero in
positive degree and $\lim_m D_{q,m}$ in degree zero. The tuple
resolution consequently computes the native cohomology by the
cosimplicial ring of these limits. This argument retains all affine
test rings and all orders allowed in the original bar definition.
% END D72P1

% BEGIN D72P2
\paragraph{D72P2. Reduction along a bounded divided-power ideal.}
Work first on an affine chart of $U$ with \'etale coordinates over
$R[t_1,\ldots,t_d]$. With the first lift held fixed, the other $q$
lifts are determined by the differences of their coordinates from
the first. Any such differences in a test ideal give lifts, uniquely
by formal \'etaleness across a nilpotent ideal. Thus the envelope in
D72P1 is described by the freely adjoined divided-power differences
$u_1,\ldots,u_{qd}$ over $A$, at finite order. The divided powers
on $IA$ extend those on $I$ by flatness of $A/R$. Write $K$ for
the difference ideal. On the absolute diagonal the order filtration
is $K^{[m]}$. On the reduced diagonal it is $(IA+K)^{[m]}$ in
the same divided-power polynomial algebra. The addition and product
identities for divided powers give
\[
 K^{[m]}\subseteq(IA+K)^{[m]},\qquad
 (IA+K)^{[m+b]}\subseteq K^{[m]} .
\]
Indeed every term of total order at least $m+b$ either contains
base factors of total order at least $b$, hence vanishes, or has
difference factors of order at least $m$. The quotients are therefore
cofinal and the two inverse limits agree. The map induced by the
given reduction is the identity on the coefficients and the universal
differences. The universal description makes it commute with deleting
and repeating lifts, even if the first lift changes. It is an isomorphism
of the two cosimplicial rings.

For clarity about passage from these charts to schemes, equip the
affine tests with the covers by principal opens of $\operatorname{Spec}B$.
They induce the same opens on $\operatorname{Spec}B_0$, since $J$ is nilpotent,
and localization carries the prescribed divided powers.
Category cohomology of the structure presheaf agrees with its sheaf
cohomology for these covers. One way to see this is to derive the
inclusion of sheaves into presheaves. At a test its higher direct images
are the cohomology on the slice over that test. This is ordinary Zariski
cohomology of $\mathcal O$ on $\operatorname{Spec}B$: covers of
the test, their open intersections and their further covers are
open localizations. Equivalently the open hypercovers used to
compute its cohomology are cofinal, by refining each matching cover
by open covers beginning at the test. Higher cohomology vanishes
on this affine. In degree zero the inclusion just recovers $B$.
Hence deriving global sections gives the original bar cohomology.
An open cover of $X$ now covers the terminal sheaf of the test site:
pull it back along each test's point and refine by principal opens.
Its intersections give the test sites of the intersections of the
chosen opens. Descent (the augmented \v Cech resolution of a cover)
thus glues the chartwise quasi-isomorphisms. Choose a smooth
coordinate cover and refine the intersections by such charts;
separatedness makes intersections of affine opens affine. The
source's quasi-compact case has a finite initial cover.

We obtain the actual pullback isomorphisms in both degrees of
$\mathrm{CrystallineDeRham}$ for every smooth separated
quasi-compact lift. The source's $m$ specifies compatibility of
the base map and the structure maps; on tests the pullback is the
one just computed. All the comparisons respect the cup product,
given on the nerve by the front and back faces and multiplication.
% END D72P2

% BEGIN D72P3
\paragraph{D72P3. Comparison with the ordinary de Rham complex.}
For the absolute site of a smooth affine coordinate chart set
$E_q=\lim_m D_{q,m}$. The following double complex also identifies
the filtration and products used for abelian schemes. Keep $q+1$
lifts and add a distinguished last lift $y$. Take
\[
 K^{q,i}=E_{q+1}\,\widehat\otimes_{A,y}\Omega^i_{A/R},
 \qquad q,i\ge0 .
\]
The horizontal faces omit one of the first $q+1$ lifts and keep $y$.
For fixed $i$ the horizontal augmentation from $\Omega^i_{A/R}$
is a homotopy equivalence. Evaluation with the last ordinary lift
set equal to $y$ gives the extra codegeneracy: on tuples append
$y$ to the list preceding $y$. The alternating face identities give
the contraction; coefficients of forms are evaluated at the same
distinguished lift throughout.

The vertical differential differentiates in $y$ while the first
$q+1$ lifts are fixed, with $d u^{[j]}=u^{[j-1]}du$ for a difference
coordinate. It agrees with the ordinary de Rham differential on
forms pulled from $y$. Relative to the first lift the vertical complex
is the completed divided-power de Rham complex in $d$ freely varying
coordinates over $E_q$. In one coordinate its augmentation $E_q$ is
a homotopy equivalence: send $u^{[j]}du$ to $u^{[j+1]}$ and kill
degree zero under the homotopy. Then $dh+hd$ is the identity
minus evaluation at zero. This raises the divided-power index by
one in the integration step and involves no scalar division.
Iterate the contractions one coordinate at a time, using evaluation
in the coordinates already contracted. Both the maps and their
identities pass to the completion: a differential lowers order by
at most one and each coefficient of the resulting series is defined
at finite order. Thus each vertical augmentation from $E_q$ is a
quasi-isomorphism. The two spectral sequences of the first-quadrant
double complex compare the bar cohomology of D72P1 with ordinary
de Rham hypercohomology. Both augmentations are functorial in maps of
smooth lifts; the coordinates were used only to prove they are
quasi-isomorphisms. Multiplication of forms and the front/back
face product respect the augmentations. They glue over the covers of
D72P2, with forms from the lift going to their pullbacks at $y$.
The first Hodge step in the source is the kernel of restriction to
the full subcategory of tests with $J=0$. Such tests are equivalent
to ordinary affines over $X$: the surjection with zero kernel
identifies $B_0$ with $B$. On an affine chart the lifts at a trivial
test are uniquely determined. Restriction of the envelope complex
to these tests is diagonal evaluation $E_q\to A$ for every $q$.
The double complex maps to the constant horizontal cosimplicial ring
$A$ by diagonal evaluation in form degree zero and by zero in positive
form degree. This commutes with its horizontal faces, and the target
has no positive vertical degrees. On the horizontal augmentation the
map is precisely the projection $\Omega^\bullet_{A/R}\to A[0]$.
Consequently the global restriction, after descent on the trivial tests
as in D72P2, is the de Rham edge map to $H^1(X,\mathcal O_X)$.
The short exact sequence with kernel $\Omega^{\ge1}_{X/R}$ identifies
its kernel with the image of global closed one-forms in de Rham
cohomology. For an abelian scheme these forms are invariant and
closed. The source's \texttt{PrimitiveSubmodule} is the intersection
of this kernel with primitive classes. Invariant forms are primitive
under the group law, so it equals the usual degree-one Hodge
submodule. The source defines the homological filtration as its
annihilator in the dual of the primitive part.

There is also a representative-level check for the source's tangent
evaluation. Integrating a closed form $\omega$ in the varying lift
with the other lift fixed gives its divided-power difference
$P(x,y)$, normalized by $P(x,x)=0$. Its differential in $y$ is
$\omega(y)$. The contraction above gives uniqueness of this
normalized primitive, proving both independence of the coordinates
and the cocycle identity $P(x,z)=P(x,y)+P(y,z)$ on triples.
For lifts on a test, evaluate the bar cocycle on an arrow as the
difference $P(\text{transported source},\text{chosen target})$. This convention gives $P(x,y)$ for the source's value
on the section arrow minus the identity arrow when $x$ is the
constant lift and $y$ the given lift (equivalently check the sign on
$\omega=da$, where the result is $a(y)-a(x)$).
On a trivial thickening the lift is forced, so the cocycle vanishes
on arrows between trivial tests. Lifts exist Zariski locally by
smoothness and on each affine test by successive square-zero lifting:
at each step the obstruction to gluing is in $H^1$ of a
quasi-coherent tangent module on an affine scheme. The construction
therefore also gives a bar representative on the full category.
For $y=e+\epsilon v$ in the source's split dual numbers, the
linear term of $P(e,y)$ is $\epsilon\omega_e(v)$ and every higher
divided power of $\epsilon R$ is zero for the prescribed divided
powers, including in characteristic two. Changing the representative
by $d^0g$ while retaining its vanishing on trivial arrows makes $g$
on trivial tests a natural function, hence a global section of
$\mathcal O_X$. Its contribution is $g(y)-g(e)$. On an abelian
scheme over $R$ every global function comes from $R$, so this
contribution vanishes. This proves the claimed evaluation for every
native form representing the same class.
% END D72P3
